\section{Discussion}
\label{sec:discussion}

\paragraph{Group sizes in practice.} Measured on enumerable policies, on transformers trained from
scratch and on a pretrained model, the purely statistical optimum came out between three and six
responses per prompt. Published recipes commonly use eight to sixty-four, and
Theorem~\ref{thm:split}(iii) accounts for the difference without either side being wrong: prefill
sharing and a large rollout budget both push the optimum up, and at the prefill ratios we measure
for long prompts with short answers they push it into the tens. What the theorem removes is the
need to guess. Both inputs are measurements, one from the trainer and one from the serving stack,
and the second changes whenever the prompt-to-response length ratio does.

\paragraph{What the estimator costs.} Four -- or $2K$ -- gradient buffers instead of one, and no
extra responses. A trainer that already accumulates over microbatches gets the decomposition for
the cost of tagging each microbatch with which prompt block and which rollout sub-group it came
from. What it buys is $\Bcrit$, the efficiency $\rho$, and $G^{\star}$, continuously, on the run in
progress rather than from a sweep that has to be repeated when anything changes.

\paragraph{Limits.} The estimator family in \eqref{eq:family} covers advantages that depend on a
response's own reward and its group's success count. It does not cover token-level shaping, length
normalisation, clipping, or a non-binary reward, and the exact moment calculations lean on the
reward being binary. Clipping in particular makes the update a nonlinear function of the batch,
which breaks the linearity the split estimator relies on; it can be handled by instrumenting the
clipped update directly, but the closed forms would need redoing. The drift model is a second-order
expansion and will understate movement for steps large enough to leave that regime. The models measured here run from tabular policies through transformers of $25$k to $400$k
parameters to a pretrained model of $0.5$B, so the claim that these are the right quantities is
tested over four orders of magnitude in parameter count and not beyond; the transfer result is
demonstrated across batch size, not across model scale, where our sweeps were not sharp enough to
separate the two parameterisations.

\paragraph{Preconditioners.} All quantities must be read in the metric the update moves in, and
Section~\ref{sec:adam} shows what that takes in practice: prime the preconditioner before measuring,
and put a floor under its second moment before inverting it, or the coordinates that rarely receive
gradient take over the inner products and the sign of $\tau_b$ stops being reliable. With both in
place the law holds under Adam as well as under gradient ascent. We still treat the preconditioner
as slowly varying relative to batch-to-batch fluctuation, which is standard and which we have not
tested at batch sizes small enough to stress it.

\paragraph{What we would measure next.} Three things. Whether $\tau_w/\tau_b$ moves once responses
are long and reasoning-heavy rather than short and exact-match, since the within-prompt term should
grow with the entropy of the response distribution. How the drift target should be scheduled rather than held constant, since a tuned
constant learning rate beat the best constant drift on the held-out setting of
Section~\ref{sec:recipe}, and whether it transfers across model scale, which needs a sharper
experiment than the one run here. And whether the same decomposition explains the well-known instability of RLVR at small batch
sizes: a run below its critical batch size spends most of its drift budget on diffusion, and the
diffusive fraction is exactly what \eqref{eq:efficiency} measures.
